\documentclass[10pt,letterpaper]{article}
\usepackage{cornell-notes}

\title{Chapter 30: Tree of Minimal Length}
\author{}
\date{}

\setDocTitle{Chapter 30: Tree of Minimal Length}
\setDocSubtitle{Combinatorial Algorithms}
\setDocVersion{v1.0}
\setDocStatus{Study Companion}
\setDocOwner{Combinatorial Algorithms Cornell Notes}
\setDocAuthor{}
\setDocDate{\today}
\setDocSummary{Cornell-format study and review notes for Chapter 30: Tree of Minimal Length.}

\setCornellCollection{Combinatorial Algorithms Cornell Notes}
\setCornellUnitType{Chapter}
\setCornellUnitNumber{30}
\setCornellUnitTitle{Tree of Minimal Length}
\setCornellCompanionLabel{Study and review companion}

\hypersetup{pdftitle={Chapter 30: Tree of Minimal Length},pdfsubject={Cornell notes},pdfauthor={}}

\begin{document}
\maketitle
\begin{tcolorbox}[colback=Paper,colframe=Ink]{\small\bfseries\color{Teal}CORNELL NOTES}\par{\LARGE\bfseries Chapter 30: Tree of Minimal Length (MINSPT)}\par\medskip\textbf{Topic:} Minimum spanning tree from a dense distance matrix\hfill\textbf{Date:} \today\par\textbf{Study purpose:} Maintain nearest-tree links while growing an optimal tree.\end{tcolorbox}
\begin{tcolorbox}[colback=Teal!5,colframe=Teal,title=Essential question]Why does repeatedly selecting the lightest edge crossing from the current tree to an outside vertex produce a spanning tree of minimum total length?\end{tcolorbox}
\begin{tcolorbox}[colback=white,colframe=Mist,title=Learning objectives]\begin{itemize}\item State the cut-based greedy rule.\item Interpret the signed TREE array.\item Trace selection and nearest-parent updates.\item Prove optimality by exchange.\item Analyze dense-matrix cost and disconnected inputs.\end{itemize}\textbf{Prerequisites:} weighted graphs, trees, cuts, asymptotic complexity.\end{tcolorbox}
\section*{Cornell notes}
\begin{longtable}{@{}Q!{\color{Mist}\vrule width 1pt}A@{}}\cornellhead
\cue{Greedy construction}{Begin with $T_0=\{n\}$. At each step select a vertex $i$ outside the current vertex set whose cheapest connection to the set has minimum length, add $i$, and record that cheapest edge. After $n-1$ additions the recorded edges form a spanning tree.}
\cue{Signed TREE meaning}{TREE$(i)=-m$ means $i$ is outside and its closest known vertex inside the current tree is $m$. TREE$(i)=m>0$ means $i$ has joined through edge $(i,m)$. TREE$(n)=0$ marks the initial root and has no parent.}
\cue{Initialization}{With only vertex $n$ inside, set TREE$(i)=-n$ for $i<n$. The provisional crossing length for $i$ is then $D(i,n)$. This establishes the nearest-tree invariant before the first selection.}
\cue{Selection step}{Among all negative entries, choose $i$ minimizing $D(i,|\mathrm{TREE}(i)|)$. Negate TREE$(i)$ to make its saved closest vertex the permanent parent and to mark $i$ as inside.}
\cue{Update step}{For every remaining outside vertex $j$, compare its saved crossing length with $D(j,i)$. If the new vertex is closer, replace TREE$(j)$ by $-i$. Thus each negative entry always names the best crossing edge currently known for that vertex.}
\cue{Why no cycle appears}{Each selected edge joins one outside vertex to the existing tree. It cannot connect two already included vertices, so it adds one vertex and one edge without creating a cycle. After $n-1$ steps the result is connected and acyclic.}
\cue{Exchange proof}{Let $e$ be the lightest edge crossing the cut between the current tree vertices and the rest. If an optimal spanning tree omits $e$, adding it creates a cycle containing another edge $f$ crossing the same cut. Since $w(e)\leq w(f)$, replacing $f$ by $e$ preserves optimality. Induction makes every greedy choice safe.}
\cue{Complexity}{With a dense $n\times n$ distance matrix, each of $n-1$ iterations scans outside vertices to select and update them. The running time is $O(n^2)$ and the extra working storage is $O(n)$ beyond the input matrix.}
\cue{Missing edges}{Represent an absent edge by $+\infty$. If the smallest provisional length is still infinite, the graph is disconnected and no spanning tree exists. The routine should report failure rather than return infinite-length parent links. Negative real weights are allowed by the proof.}
\cue{Retrieval practice}{\begin{enumerate}\item What does a negative TREE entry encode?\item Which edge is selected next?\item Why does one update scan suffice?\item What edge is exchanged in the proof?\item How is disconnection detected?\end{enumerate}}
\end{longtable}
\begin{tcolorbox}[colback=Ink!4,colframe=Ink,title=Cornell summary]
MINSPT grows a minimum spanning tree from one starting vertex. For each outside vertex, a negative TREE entry stores its nearest current-tree neighbor; the lightest of those provisional crossing edges is chosen, its sign is flipped, and the new vertex is used to improve the remaining links. Every addition connects one new vertex and cannot form a cycle. The cut exchange argument proves the choice is compatible with some optimum at every step. Dense-matrix selection and updates take $O(n^2)$ time and $O(n)$ auxiliary storage.
\end{tcolorbox}
\end{document}
